% Generated by scripts/assemble.py; corrections belong in source/corrections.json.
\examhead{9}{北京工业大学 2025 年数学分析试卷}{15}
% p016-b09; source PDF page 16
\par\Needspace{6\baselineskip}
\blockanchor{p016-b09}
\question{1}(15 分) 设 p > 1, 证明: 数列  $\left\{1 + \frac{1}{2^{p}} + \cdots + \frac{1}{n^{p}}\right\}$  收敛, 且其极限 a 满足\par
\nopagebreak[4]
% p016-b10; source PDF page 16
\blockanchor{p016-b10}
\displayformula{1 + \frac {1}{2 ^ {p - 1} (p - 1)} \leq a \leq \frac {p}{p - 1}.}
% p016-b11; source PDF page 16
\blockanchor{p016-b11}
\question{2}(15 分) 证明: 方程组 $\frac{5}{x - 1} + \frac{7}{x - 2} + \frac{16}{x - 3} = 0$ 在 (1,2) 与 (2,3) 内各有一实根.\par
% p016-b12; source PDF page 16
\par\Needspace{6\baselineskip}
\blockanchor{p016-b12}
\question{3}(15 分) 设函数 $f(x)$ 在 $[0,1]$ 上可导, $f(0) = 0$, $f(1) = 2$. 证明: 在 $(0,1)$ 上存在两个 $\xi_1$ 与 $\xi_2$, 使得\par
\nopagebreak[4]
% p016-b13; source PDF page 16
\blockanchor{p016-b13}
\displayformula{\frac {1}{f ^ {\prime} (\xi_ {1})} + \frac {1}{f ^ {\prime} (\xi_ {2})} = 1.}
% p016-b14; source PDF page 16
\blockanchor{p016-b14}
\question{4}(15 分) 证明: 若函数 $f(x)$ 在 $[a, +\infty)$ 上连续, 且 $\lim_{x \to +\infty} [bx - f(x)] = 0$, 其中 $b$ 为非零实数. 则 $f(x)$ 在 $[a, +\infty)$ 上一致连续.\par
% p016-b15; source PDF page 16
\blockanchor{p016-b15}
\question{5}(15 分) 证明: 存在常数 $A < 1$, 对任意的正整数 $n > 1$, 都有 $\sqrt{1} + \sqrt{2} + \cdots + \sqrt{n} < An^{\frac{3}{2}}$.\par
% p016-b16; source PDF page 16
\blockanchor{p016-b16}
\question{6}(15 分) 判别函数列 $\{nx(1 - x)^n\}$ 在 $[0, 1]$ 上的一致收敛性.\par
% p016-b17; source PDF page 16
\blockanchor{p016-b17}
\question{7}(15 分) 求幂级数 $\sum_{n=1}^{\infty} n^{2} x^{n+1}$ 的收敛域与和函数.\par
% p016-b18; source PDF page 16
\blockanchor{p016-b18}
\question{8}(15 分) 求函数 $f(x, y, z) = x - 2y + 2z$ 在条件 $x^2 + y^2 + z^2 = 1$ 的约束下的最大值和最小值.\par
% p016-b19; source PDF page 16
\blockanchor{p016-b19}
\question{9}(15 分) 求函数 $f(x) = \int_{0}^{+\infty} \frac{\ln(1 + y^3)}{y^x} \, \mathrm{d}y$ 的定义域.\par
% p016-b20; source PDF page 16
\par\Needspace{6\baselineskip}
\blockanchor{p016-b20}
\question{10}(15 分) 记\par
\nopagebreak[4]
% p016-b21; source PDF page 16
\blockanchor{p016-b21}
\displayformula{F (t) = \iiint_ {V} f \left(x ^ {2} + y ^ {2} + z ^ {2}\right) \mathrm{d} x \mathrm{d} y \mathrm{d} z.}
% p016-b22; source PDF page 16
\blockanchor{p016-b22}
其中 $V = \{(x,y,z)\mid x^2 +y^2 +z^2\leq t^2\}$$(t > 0)$ ，若 $f$ 可微，求 $F^{\prime}(t)$\par
